% Generated draft: requires mathematical and source review.
\documentclass[11pt]{article}
\usepackage[margin=25mm]{geometry}
\usepackage{fontspec}
\setmainfont{Latin Modern Roman}
\usepackage{amsmath,amssymb,mathtools,unicode-math}
\setmathfont{Latin Modern Math}
\usepackage{enumitem,needspace,xurl,hyperref}
\setlength{\parindent}{0pt}
\setlength{\parskip}{5pt}
\setlength{\emergencystretch}{3em}
\newcommand{\N}{\mathbb{N}}
\newcommand{\Z}{\mathbb{Z}}
\newcommand{\Q}{\mathbb{Q}}
\newcommand{\R}{\mathbb{R}}
\newcommand{\C}{\mathbb{C}}
\newcommand{\F}{\mathbb{F}}
\newcommand{\A}{\mathbb{A}}
\newcommand{\PP}{\mathbb{P}}
\newcommand{\E}{\mathbb{E}}
\newcommand{\eps}{\varepsilon}
\newcommand{\dd}{\,\mathrm{d}}
\begin{document}
\section*{UnsolvedMath ID: 30003662}
\textbf{Source problem code:} catalogue entry
\subsection{Definitions and mathematical statement}
The supplied problem asks: for natural classes of provably convergent sequences, which uniform rates, moduli, or other quantitative forms of convergence can be extracted? It is a research programme whose precise hypotheses and desired conclusions remain unspecified in the supplied record.

The following gives a conditional framework for interpreting some of the requested quantitative notions; it does not impose a domain on the original problem. Suppose that a class \(\mathcal C\) consists of sequences \(x=(x_n)_{n\in\N}\) in a metric space \((X,d)\), with \(\N=\{0,1,2,\ldots\}\) as an indexing convention for this framework. If a limit \(\ell_x\in X\) is specified, a rate of convergence for \(x\) is a function \(r_x:\N\to\N\) satisfying
\[
\forall k\in\N\;\forall n\ge r_x(k),\qquad d(x_n,\ell_x)\le 2^{-k}.
\]
A Cauchy modulus for \(x\) is a function \(c_x:\N\to\N\) satisfying
\[
\forall k\in\N\;\forall m,n\ge c_x(k),\qquad d(x_m,x_n)\le 2^{-k}.
\]
A uniform Cauchy modulus for \(\mathcal C\) is a single function \(c\) satisfying this condition for every \(x\in\mathcal C\). A uniform bound on metastability for \(\mathcal C\) is a functional \(\Phi:\N\times\N^{\N}\to\N\) such that
\[
\forall x\in\mathcal C\;\forall k\in\N\;\forall g:\N\to\N\;
\exists N\le\Phi(k,g)\;
\forall i,j\in[N,N+g(N)]\cap\N,\qquad
 d(x_i,x_j)\le 2^{-k}.
\]
In these definitions, uniformity means independence from the particular sequence. If dependence on additional parameters, bounds, or proofs is allowed, those inputs must instead be specified explicitly.

\textbf{Statement to be made precise.} Specify a formal theory \(T\), a class \(\mathcal C\), a notion of convergence, and the representation and hypotheses under which the relevant convergence assertions are provable in \(T\). Determine which quantitative convergence data can be extracted from such proofs, with what effectiveness and with what dependence on the sequence, proof, and other permitted parameters. Potential targets include rates of convergence, Cauchy moduli, bounds on metastability, and other quantitative data.

\textbf{Unresolved qualifications.} The record specifies neither \(T\), \(\mathcal C\), the ambient space or convergence notion, nor the permitted parameter dependence and meaning of extraction. It also does not distinguish individual provability for each sequence from a proof of a universally quantified convergence assertion for a class. Thus the metric definitions above illustrate possible targets, while the supplied evidence does not determine a single fully quantified conjecture.
\subsection{Short English statement}
The problem asks what quantitative information can be extracted from proofs that sequences converge, uniformly over suitable natural classes. Possible information includes convergence rates, Cauchy moduli, and metastability bounds. The formal theory, classes, convergence notion, and permitted dependencies must be specified before this becomes a precise individual problem. The supplied status is \emph{partially solved}; this is a curation label, not an independently verified mathematical conclusion.
\subsection{Sources}
\begin{itemize}
\item Supplied local UnsolvedMath registry record, permanent ID \(30003662\), title \emph{Uniform Quantitative Rates for Provably Convergent Sequences}, source problem code \texttt{catalogue entry}, in the \emph{Oberwolfach Reports Open Problems} collection. The record supplies proposed year 2017 and no individual proposer.
\item Supplied report citation: \emph{Mathematical Logic: Proof Theory, Constructive Mathematics} (2018), Oberwolfach Reports, pp.\ 3121--3183, DOI \texttt{10.4171/owr/2017/53}. \url{https://doi.org/10.4171/owr/2017/53}. The supplied background describes Oberwolfach Reports as documenting mathematical workshops and their open-problem discussions.
\item Supplied literature source: OWR Report 53/2017. \url{https://ems.press/content/serial-article-files/46716}.
\item Supplied literature source: Howard--Towsner abstract. \url{https://mca2017.mcofamericas.org/sites/mca2017/files/abstracts/ibm-ht.pdf}. No further bibliographic details or attribution of the problem to these authors are supplied.
\item Literature assessment dated exactly \texttt{2026-08-21}: supplied status \emph{partially solved}. The supplied assessment states that proof-theoretic work establishes uniform metastability and uniform jump bounds for broad classes of provably convergent sequences, but that no universal effective rate answers the unrestricted question. This is supplied dated curation derived from the OWR material, not independently verified here; the time-sensitive status and cited claims require independent verification.
\item The registry labels the statement \texttt{corrected\_verified}, explaining that mathematical content was preserved while notation, typography, and sentence structure were repaired for publication. This records the supplied editorial assessment and does not establish mathematical correctness or independent source verification for this draft.
\end{itemize}
\end{document}
